\chapter{Discussion}
\label{chap:discussion}

\section{Answers to the research questions}

\textbf{RQ1: research landscape.} The selected two-venue corpus increases from
2,240 papers in 2021 to 4,042 in 2025, with a dip in 2022. The embedding
classifier assigns topics to 85.00\% of papers. Within labeled papers, the
multimodal share rises in both venue series. These are measured changes in the
selected proceedings and model-defined topic assignments, not evidence of
field-wide growth or independently verified topic shifts.

\textbf{RQ2: collaboration structure.} The researcher graph is sparse on average
(mean distinct degree 5.00), has a large component containing 64.5\% of resolved
researchers, and includes many small components and isolates. Organization
collaboration includes substantial within-academia support and 3,533 distinct
academia--industry pair-paper supports. The structure supports candidate maps for
further review, not claims about current partnerships or employment.

\textbf{RQ3: bridges and citation attention.} Bo An is the highest-ranked
researcher under the configured approximate betweenness measure, while the
Informer paper leads in in-corpus citation in-degree with 51 incoming edges.
Both measures depend on a restricted graph and have different limitations:
identity-policy changes materially alter bridge rankings, and citation counts
exclude external targets and do not normalize paper age.

\textbf{RQ4: reliability.} Structural and database integrity checks pass, but
these validate consistency rather than meaning. The v4 global audit passes its
numeric thresholds only under AI judgments; the historical v3 audit failed its
numeric checks. Topic agreement is modest against AI-assigned reference labels,
affiliation coverage misses its target, and bridge rankings are sensitive to
identity and evidence choices. Verified course accuracy remains unestablished.

\section{Methodological strengths}

The construction has a clear corpus boundary and source hierarchy. DOI-derived
matching, title similarity and byline-count checks restrict work-level joins.
Author-position compatibility and global ORCID conflict handling prevent
name-only merging. Paper-specific affiliation triples preserve the evidence
needed for later filtering. Accepted edges retain deterministic IDs, source
references, method versions, and graph versions. Separate edge classes keep
observed publication facts apart from inferred topics and derived social links.

The analysis also records its projection and counting rules. Collaboration
weights retain fractional contributions, citation direction is explicit, topic
shares use fractional allocation, and all plots preserve the full graph
denominators or state when they show a bounded selection. These practices make
the results easier to reproduce and reduce ambiguity in interpretation.

\section{Limitations and threats to validity}

\subsection{Coverage and metadata}

Eleven roster papers remain unmatched to OpenAlex and 15,032 matched byline
positions are withheld. Affiliation evidence covers 73.71\% of matched
positions, below the 80\% target. Missing positions may not be random across
venues, years, or institutions, so the observed researcher and organization
networks can be a selective view. Seven historical returned-DOI disagreements
remain unresolved in v4; matching titles/bylines do not establish source-version
equivalence or validate the associated affiliations. The graph covers AAAI and IJCAI only; it cannot
support conclusions about all AI publication venues.

\subsection{Researcher identity and derived edges}

The v4 correction resolves a concrete observed-ORCID omission by withholding
provider IDs with conflicting ORCID evidence. It does not provide independent
identity verification for every remaining OpenAlex author ID. The six global
v4 audit uncertainties are identity-related, and coauthorship, organization
collaboration, and researcher centrality inherit identity errors from their
inputs. Since identity candidates change the largest component and bridge
rankings, stable IDs in a provider are not equivalent to guaranteed semantic
identity.

\subsection{Topic labels and temporal interpretation}

Topic assignment is inferred from titles and abstracts under one fixed model,
taxonomy, and threshold set. Coverage is not accuracy. The historical evaluation
uses exposed test papers and AI-generated references, with embedding micro F1
0.4667 and phrase-baseline F1 0.2301. The v4 audit creates no new topic labels.
Consequently topic profiles, organization portfolios, and year-to-year shifts
should be treated as exploratory. Abstained papers may have a different topic
mix from labeled papers.

\subsection{Evaluation dependence}

AI reviewers used the same general model family and construction sources. Fresh
contexts avoid shared answer exposure but do not establish reviewer or model
independence. The prospective audit corrected a graph after an earlier audit
found a resolver issue; the new source snapshot and sample history are therefore
not equivalent to a pristine external test. The Wilson interval covers sampling
uncertainty conditional on judgments only. It does not cover evaluator errors,
source overlap, systematic model bias, or identity mapping uncertainty. Human
validation was not performed.

\subsection{Network measures and business use}

Betweenness depends on edge weights, the inverse-strength distance rule, the
projection, and the approximate pivot set. A high rank is a structural signal
for inspection, not a person's quality or causal influence. Citation in-degree
uses only within-corpus edges and gives older papers more exposure time. Sector
and country labels are institutional metadata, not attributes of individual
researchers. None of the reported network measures implies hiring value,
commercial readiness, or a current institutional partnership.

\section{Practical implications}

For a research-program analyst, the graph can help locate repeated
co-publication supports, examine bridge candidates, and compare topic portfolios
within the selected venues. The appropriate next action is source inspection:
check supporting papers, author positions, institution strings, and topic
definitions before using a shortlist. The report does not rank organizations or
people as business recommendations. It presents an evidence-linked starting
point and makes the uncertainty visible alongside the result.

\section{Future work}

Future work should prioritize independent, source-backed annotation for a
properly frozen sample; identity resolution with person-specific evidence;
additional source coverage for unmatched or affiliation-poor papers; and a
topic evaluation with an untouched human reference set. Recall requires an
evaluation frame that includes plausible missing edges, not only accepted graph
claims. Temporal citation normalization and documented sensitivity analysis
would improve comparisons across years. Any changes to the production graph
should create a new version and a new audit rather than overwrite v4.
